%% SPDX-License-Identifier: GPL-2.0-or-later
%% Copyright (C) 2026 the DoomTeX authors
%%%%%%%%%%%%%%%%%%%%%%%%%%%%%%%%%%%%%%%%%%%%%%%%%%%%%%%%%%%%%%%%%%%%%%%%%%%%%%
%%% wad.tex -- read a real Doom WAD in pure pdfTeX.
%%%
%%% \pdffiledump returns the hex bytes of any file region, expandably and with
%%% no shell escape, so TeX can parse the actual binary format: the 12-byte
%%% header, the lump directory, and the map's VERTEXES / LINEDEFS / SIDEDEFS /
%%% SECTORS / SEGS / SSECTORS / NODES / THINGS structures.  PLAYPAL is copied
%%% out verbatim, because a Doom palette and a PDF /Indexed colour space are
%%% the same 768 bytes.
%%%
%%% Input by doom.tex when \DTWAD is defined, and it fills exactly the arrays
%%% the built-in map does -- through the same \dtAdd... macros -- so nothing
%%% downstream knows which of the two it is drawing.  In file order: integer
%%% decoding and the lump directory, COLORMAP, sprites, the status bar's
%%% pictures, wall textures and their masked middles, flats, and the loader
%%% that calls all of them.
%%%
%%%   \dtLoadWAD{doom1.wad}{E1M1}
%%%%%%%%%%%%%%%%%%%%%%%%%%%%%%%%%%%%%%%%%%%%%%%%%%%%%%%%%%%%%%%%%%%%%%%%%%%%%%
\catcode`\@=11

%%% -------------------------------------------------------- hex digit values
\dt@q=0
\dtWhile{\dt@q<16}{%
  \expandafter\xdef\csname dtHV\dt@hd{\the\dt@q}\endcsname{\the\dt@q}%
  \advance\dt@q by1\relax }
%% accept lowercase too, in case a distribution differs
\expandafter\gdef\csname dtHVa\endcsname{10}
\expandafter\gdef\csname dtHVb\endcsname{11}
\expandafter\gdef\csname dtHVc\endcsname{12}
\expandafter\gdef\csname dtHVd\endcsname{13}
\expandafter\gdef\csname dtHVe\endcsname{14}
\expandafter\gdef\csname dtHVf\endcsname{15}

%%% ------------------------------------------- little-endian integer decoding
%%% NB: names are letters only.  \dtU32 would parse as \dtU followed by the
%%% characters 32, because digits are not letters -- so \dtU16 and \dtU32
%%% would silently define the same macro with conflicting parameter texts.
\def\dtB#1#2{\csname dtHV#1\endcsname*16+\csname dtHV#2\endcsname}
\def\dtUW#1#2#3#4{\number\numexpr\dtB{#1}{#2}+256*(\dtB{#3}{#4})\relax}
\def\dtSW#1#2#3#4{\number\numexpr\dtUW{#1}{#2}{#3}{#4}%
  -\ifnum\dtUW{#1}{#2}{#3}{#4}>32767 65536\else0\fi\relax}
\def\dtUL#1#2#3#4#5#6#7#8{\number\numexpr\dtB{#1}{#2}+256*(\dtB{#3}{#4})%
  +65536*(\dtB{#5}{#6})+16777216*(\dtB{#7}{#8})\relax}

%%% \dtRead<n>{count}{offset} : read n bytes and leave the decoded value
\def\dtWADFILE{}
\def\dtRd#1#2{\csname dtRd@#1\endcsname{\pdffiledump offset #2 length #1
  {\dtWADFILE}}}
\expandafter\def\csname dtRd@2\endcsname#1{\expandafter\dtSW#1}
\expandafter\def\csname dtRd@2u\endcsname#1{\expandafter\dtUW#1}
\expandafter\def\csname dtRd@4\endcsname#1{\expandafter\dtUL#1}
\def\dtS#1{\csname dtRd@2\endcsname{\pdffiledump offset #1 length 2
  {\dtWADFILE}}}
\def\dtUu#1{\csname dtRd@2u\endcsname{\pdffiledump offset #1 length 2
  {\dtWADFILE}}}
\def\dtL#1{\csname dtRd@4\endcsname{\pdffiledump offset #1 length 4
  {\dtWADFILE}}}

%%% ------------------------------------------------- lump name -> hex pattern
%%% Lump names are 8 bytes, NUL padded.  Comparing them as hex strings avoids
%%% any catcode trouble with the raw bytes.
\newcount\dt@nc
\def\dtNameHex#1{\gdef\dt@target{}\dt@nc=0 \dt@nhA#1\relax
  \dtWhile{\dt@nc<8}{\xdef\dt@target{\dt@target00}\advance\dt@nc by1\relax }}
\def\dt@nhA#1{\ifx\relax#1\else
  \xdef\dt@target{\dt@target\hx{`#1}}\advance\dt@nc by1
  \expandafter\dt@nhA\fi}

%%% --------------------------------------------------------- directory lookup
\newcount\dtNUML \newcount\dtDIR \newcount\dtENT
\newcount\dtLPOS \newcount\dtLSIZE \newcount\dtFOUND

%% \dtFindLump{NAME}{startentry}: sets \dtFOUND (entry index or -1),
%% \dtLPOS, \dtLSIZE
\def\dtFindLump#1#2{%
  \dtNameHex{#1}%
  \dtFOUND=-1 \dtENT=#2\relax
  \dtWhile{\dtENT<\dtNUML}{%
    \ifnum\dtFOUND<0
      \edef\dt@nm{\pdffiledump offset
        \number\numexpr\dtDIR+16*\dtENT+8\relax\space length 8 {\dtWADFILE}}%
      %% \ifx would compare catcodes too: \pdffiledump yields catcode-12
      %% characters while the hex table produces catcode-11 letters for A-F.
      %% \pdfstrcmp compares character codes only.
      \ifnum\pdfstrcmp{\dt@nm}{\dt@target}=0
        \dtFOUND=\dtENT
        \dtLPOS=\dtL{\number\numexpr\dtDIR+16*\dtENT\relax}\relax
        \dtLSIZE=\dtL{\number\numexpr\dtDIR+16*\dtENT+4\relax}\relax
      \fi
    \fi
    \advance\dtENT by1\relax }}

%%% ==============================================================  COLORMAP
%%% 34 tables of 256 bytes: COLORMAP[light][colour] is the palette index to
%%% draw instead of `colour' at that light level.  The renderer's own surfaces
%%% are shaded by adding 0..13 to a ramp base, which works only because those
%%% bases are hand-picked rows that really do darken -- see the PLAYPAL note
%%% in INTERNALS.md.  Real artwork has no such freedom: a sprite pixel is whatever
%%% colour the artist used, so shading it needs Doom's own table.
%%%
%%% Keyed by the two hex characters of the source byte, because that is
%%% exactly the form sprite pixels arrive in.  Nothing is decoded on the way
%%% in or on the way out.
\newcount\dtCML \newcount\dtCMJ
%% Without a COLORMAP the walls would still dereference dtCM<light>@<byte>,
%% and an undefined csname inside a nested \csname is a "Missing \endcsname"
%% rather than a wrong pixel.  One identity row keeps that lookup total and
%% \dtCMN pins the light to it, so a PWAD with no COLORMAP draws its textures
%% at full brightness instead of failing to compile.
\def\dt@identitycm{%
  \dtCMJ=0
  \dtWhile{\dtCMJ<256}{%
    \expandafter\xdef\csname dtCM0@\hx{\the\dtCMJ}\endcsname{\hx{\the\dtCMJ}}%
    \advance\dtCMJ by1\relax }%
  \global\dtCMN=1 }
\def\dtLoadColormap{%
  \dtFindLump{COLORMAP}{0}%
  \ifnum\dtFOUND<0
    \message{[wad] no COLORMAP -- sprites will draw unshaded.}%
    \dt@identitycm
  \else
    %% Thirty-three tables, not thirty-two.  The lump holds thirty-four: 0..31
    %% are the light levels, which is what \dtCMN counts and what every clamp
    %% in the renderer keeps inside, and table 32 is the inverse map Doom
    %% draws the whole view through while the player is invulnerable.  It is
    %% reached only through \dtFIXCM, never by a light level.
    \dtCML=0
    \dtWhile{\dtCML<33}{%
      \edef\dt@cm{\pdffiledump offset
        \number\numexpr\dtLPOS+256*\dtCML\relax\space length 256
        {\dtWADFILE}}%
      \dtCMJ=0
      \expandafter\dt@cmscan\dt@cm
      \advance\dtCML by1\relax }%
    \global\let\dt@srPx\dt@srCM
    \global\dtCMN=32
    \message{[wad] COLORMAP: 32 light levels}%
  \fi}
%% Consumes a row two characters at a time.  \hx builds the key in the same
%% uppercase hex \pdffiledump produces, so a lookup from sprite data hits.
\def\dt@cmscan#1#2{%
  \expandafter\xdef\csname dtCM\the\dtCML @\hx{\the\dtCMJ}\endcsname{#1#2}%
  \advance\dtCMJ by1
  \ifnum\dtCMJ<256 \expandafter\dt@cmscan\fi}

%%% ================================================================  SPRITES
%%% Sprite lumps live between the S_START and S_END markers and are named
%%% <four characters><frame><rotation>, with an optional second frame and
%%% rotation in bytes 6 and 7 when one picture serves two frames mirrored.
%%% Matching only the first six bytes covers both spellings: POSSA1 is padded
%%% with NULs past byte 6, and BOS2A1C1 carries its second frame there.
%%%
%%% Rotation 1 is the view of a thing facing you, so using it for every angle
%%% gives sprites that always turn to the camera.  Rotation 0 means the
%%% picture is the same from all sides, which is how every pickup is stored.
\newcount\dtQS \newcount\dtQE \newcount\dtQJ
\newcount\dtQN \newcount\dtQG
\newcount\dtQW \newcount\dtQH \newcount\dtQC
\newcount\dtQP \newcount\dtQT \newcount\dtQL \newcount\dtQK
\newcount\dtQD
\dtQN=0

%% one unsigned byte
\expandafter\def\csname dtRd@1\endcsname#1{%
  \number\numexpr\expandafter\dtB#1\relax}
\def\dtBy#1{\csname dtRd@1\endcsname{\pdffiledump offset #1 length 1
  {\dtWADFILE}}}

%% like \dtNameHex but with no padding, for a partial-name comparison
\def\dtNameHexRaw#1{\gdef\dt@target{}\dt@nc=0 \dt@nhA#1\relax}

%% \dtWantSprite{KEY} -- this map needs the picture for KEY (four characters
%% and a frame letter).  Collected first and resolved in a single pass,
%% because \dtFindLump scans the directory from a given entry and a WAD this
%% size has a few thousand of them.
%% The name and the frame letter are what identify a picture; the rotation
%% is read back off whichever lumps match, because a thing needs all eight
%% and four of them are stored inside the other four.
\def\dtWantSprite#1{%
  \ifcsname dtWX#1\endcsname\else
    \expandafter\gdef\csname dtWX#1\endcsname{1}%
    \expandafter\xdef\csname dtWK\the\dtQN\endcsname{#1}%
    \dt@splitkey#1\relax
    \global\advance\dtQN by1
  \fi}
%% four characters of name, then the frame letter
\def\dt@splitkey#1#2#3#4#5\relax{%
  \dtNameHexRaw{#1#2#3#4}%
  \expandafter\xdef\csname dtWN\the\dtQN\endcsname{\dt@target}%
  \dtNameHexRaw{#5}%
  \expandafter\xdef\csname dtWR\the\dtQN\endcsname{\dt@target}}

%% every frame of an animated type.  \dt@wbase/\dt@wty rather than the
%% \dt@base/\dt@ty doom.tex uses for the same job at draw time: these run at
%% load and those run per frame, so they do not overlap today -- and a name
%% shared across that boundary is exactly what cost a day on \dt@tr.
\def\dt@wantall#1{%
  \edef\dt@wbase{\dt@dropf#1\relax}%
  %% expanded to literal letters first, for the reason \dt@animframe records
  \edef\dt@wseq{\csname dtMA\dt@wty\endcsname}%
  \expandafter\dt@wantseq\dt@wseq\relax}
\def\dt@wantseq#1{\ifx\relax#1\else
  \edef\dt@next{\noexpand\dtWantSprite{\dt@wbase#1}}\dt@next
  \expandafter\dt@wantseq\fi}

%% ...and every frame of every state a monster can be in.  Five sequences
%% rather than one, walked with the same \dt@wantseq: the state number is the
%% only thing that changes, so this is a loop over dtOa<type>@0..4 handing
%% each to the walker that already exists.
%%
%% The count register is \dtQK, which the sprite loader owns and is not
%% holding at load time.  \dtJ is NOT free -- it is the bound of the thing
%% loop this is called from, and INTERNALS section 6 records what borrowing one
%% of those costs.
\def\dt@wantstates#1{%
  \edef\dt@wbase{\dt@dropf#1\relax}%
  \dtQK=0
  \dtWhile{\dtQK<5}{%
    \ifcsname dtOa\dt@wty @\the\dtQK\endcsname
      \edef\dt@wseq{\csname dtOa\dt@wty @\the\dtQK\endcsname}%
      \expandafter\dt@wantseq\dt@wseq\relax
    \fi
    \advance\dtQK by1\relax }}

%% ...and every frame of the MISSILE a thing puts into the world, which is
%% the same argument a second level up and the first one where the pictures
%% belong to something that is not on the map at all.  A fireball is spawned
%% by an imp, so the imp is what has to be seen for BAL1 to be wanted; a
%% rocket is spawned by the player, so what has to be seen is the launcher
%% lying on the floor.  dtJm answers both.
%%
%% Two states, not five, and the sprite NAME comes out of dtJk per state
%% rather than off the key: plasma flies as PLSS and explodes as PLSE, and
%% wanting PLSSC..PLSSG instead would read five lumps that do not exist and
%% drop the explosion to a billboard.
\def\dt@wantmissile#1{%
  \dtQK=0
  \dtWhile{\dtQK<2}{%
    \ifcsname dtJa#1@\the\dtQK\endcsname
      \edef\dt@wbase{\csname dtJk#1@\the\dtQK\endcsname}%
      \edef\dt@wseq{\csname dtJa#1@\the\dtQK\endcsname}%
      \expandafter\dt@wantseq\dt@wseq\relax
    \fi
    \advance\dtQK by1\relax }}

\def\dtLoadSprites{%
  \ifnum\dtQN>0
    \dtFindLump{S_START}{0}%
    \ifnum\dtFOUND<0
      \message{[wad] no S_START -- keeping placeholder sprites.}%
    \else
      \dtQS=\numexpr\dtFOUND+1\relax
      \dtFindLump{S_END}{\the\dtQS}%
      \ifnum\dtFOUND<0 \dtQE=\dtNUML \else \dtQE=\dtFOUND \fi
      \dtQG=0 \dtQJ=\dtQS
      \dtWhile{\dtQJ<\dtQE}{%
        \edef\dt@nm{\pdffiledump offset
          \number\numexpr\dtDIR+16*\dtQJ+8\relax\space length 8
          {\dtWADFILE}}%
        \expandafter\dt@nmsplit\dt@nm
        \dtQK=0
        \dtWhile{\dtQK<\dtQN}{%
          \ifnum\pdfstrcmp{\dt@nA}{\csname dtWN\the\dtQK\endcsname}=0
            \dt@spmatch
          \fi
          \advance\dtQK by1\relax }%
        \advance\dtQJ by1\relax }%
      \message{[wad] sprites: \the\dtQG\space pictures for
        \the\dtQN\space things}%
    \fi
  \fi}
%% A lump name is eight bytes: four of name, then frame and rotation, then a
%% second frame and rotation when one picture serves two views mirrored.  One
%% dump, split into those three pieces, is cheaper than three dumps.
\def\dt@nmsplit#1#2#3#4#5#6#7#8#9{%
  \edef\dt@nA{#1#2#3#4#5#6#7#8}\dt@nmsplitB{#9}}
\def\dt@nmsplitB#1#2#3#4#5#6#7#8{%
  \edef\dt@nB{#1#2#3#4}\edef\dt@nC{#5#6#7#8}}
\def\dt@spmatch{%
  \expandafter\dt@sprot\dt@nB{0}%
  \expandafter\dt@sprot\dt@nC{1}}
%% #1#2 is the frame byte and #3#4 the rotation byte.  Rotation characters are
%% '0' to '8', hex 30 to 38, so #4 IS the rotation digit and needs no
%% decoding.  #5 says whether this slot is the mirrored one.
\def\dt@sprot#1#2#3#4#5{%
  \ifnum\pdfstrcmp{#1#2}{\csname dtWR\the\dtQK\endcsname}=0
    \dt@sptake{#4}{#5}%
  \fi}
\def\dt@sptake#1#2{%
  \advance\dtQG by1
  \dtLPOS=\dtL{\number\numexpr\dtDIR+16*\dtQJ\relax}\relax
  \edef\dt@next{\noexpand\dt@readpic
    {\csname dtWK\the\dtQK\endcsname#1}{#2}}%
  \dt@next}

%%% Doom's picture format: width, height, leftoffset and topoffset as int16,
%%% then one int32 column offset per column, then per column a chain of posts
%%% -- topdelta, length, a pad byte, `length' pixel bytes, a pad byte --
%%% ending at topdelta 255.
%%%
%%% A post's pixels are ONE \pdffiledump and are stored exactly as they come
%%% out of the file.  Reading this a byte at a time would be one file dump per
%%% pixel, which is the difference between a fraction of a second and minutes.
\def\dt@readpic#1#2{%
  \ifnum#2=1 \expandafter\gdef\csname dtPf#1\endcsname{1}\fi
  \dtQW=\dtS{\the\dtLPOS}\relax
  \dtQH=\dtS{\number\numexpr\dtLPOS+2\relax}\relax
  \expandafter\xdef\csname dtPw#1\endcsname{\the\dtQW}%
  \expandafter\xdef\csname dtPh#1\endcsname{\the\dtQH}%
  \expandafter\xdef\csname dtPl#1\endcsname
    {\dtS{\number\numexpr\dtLPOS+4\relax}}%
  \expandafter\xdef\csname dtPt#1\endcsname
    {\dtS{\number\numexpr\dtLPOS+6\relax}}%
  \dtQC=0
  \dtWhile{\dtQC<\dtQW}{%
    \dtQP=\numexpr\dtLPOS+\dtL{\number\numexpr\dtLPOS+8+4*\dtQC\relax}\relax
    \gdef\dt@post{}%
    \dtQT=0
    \dt@readposts
    \expandafter\xdef\csname dtQ#1@\the\dtQC\endcsname{\dt@post{-1}{0}{}}%
    \advance\dtQC by1\relax }}

%% \dtQT doubles as the post counter: a column of a real picture has one or
%% two posts, and the bound is only there so that a bad column offset cannot
%% spin forever.  The recursion is not tail recursion, but it is bounded by
%% that same 64 and costs nothing.
\def\dt@readposts{%
  \dtQD=\dtBy{\the\dtQP}\relax
  \ifnum\dtQD<255
    \advance\dtQT by1
    \dtQL=\dtBy{\number\numexpr\dtQP+1\relax}\relax
    \ifnum\dtQL>0
      \xdef\dt@post{\dt@post{\the\dtQD}{\the\dtQL}%
        {\pdffiledump offset \number\numexpr\dtQP+3\relax\space
         length \the\dtQL\space{\dtWADFILE}}}%
    \fi
    \dtQP=\numexpr\dtQP+4+\dtQL\relax
    \ifnum\dtQT<64 \dt@readposts \fi
  \fi}

%%% Pictures that are not sprites: no S_START range to narrow the search to,
%%% and no frame or rotation in the name.  The status bar needs some seventy
%%% of them -- digits in two fonts, a percent sign, the faces, the arms panel
%%% and the keys -- and finding each with a \dtFindLump of its own would be
%%% seventy scans of a 2928-entry directory, at one file dump per entry:
%%% some 200000 dumps against the 6000 the entire map load costs.  So this is
%%% the \dt@wantpatches shape again, and for exactly the same reason: mark
%%% the names, then walk the directory ONCE.
%%%
%%% dtGq<namehex> holds the name as it was asked for, so the pass has a key
%%% to store the picture under -- it is reading hex out of the file, not the
%%% ASCII the caller used.
\newcount\dtGN \dtGN=0
\def\dtWantGfx#1{%
  \dtNameHex{#1}%
  \expandafter\gdef\csname dtGq\dt@target\endcsname{#1}}
\def\dtLoadGfxAll{%
  \dtQJ=0
  \dtWhile{\dtQJ<\dtNUML}{%
    \dtReadName{\number\numexpr\dtDIR+16*\dtQJ+8\relax}%
    \ifcsname dtGq\dt@tnm\endcsname
      \dtLPOS=\dtL{\number\numexpr\dtDIR+16*\dtQJ\relax}\relax
      %% \dt@readpic's key must arrive as literal characters: a macro's
      %% arguments are not expanded while being grabbed, so the \csname
      %% itself would become the key and every lookup below would miss.
      \edef\dt@next{\noexpand\dt@readpic{\csname dtGq\dt@tnm\endcsname}{0}}%
      \dt@next
      \advance\dtGN by1
    \fi
    \advance\dtQJ by1\relax }}

%%% What the bar is made of.  Ten tall digits, a percent sign, five pain
%%% levels of three glances each, and the corpse face.  Written as loops
%%% rather than as thirty-one \dtWantGfx lines because the names really are
%%% derived -- and built through an \edef with the macro \noexpand'd, since
%%% \dtWantGfx walks its argument character by character and would otherwise
%%% be handed \the itself.
\def\dt@wanthud{%
  \dtWantGfx{STBAR}\dtWantGfx{STTPRCNT}\dtWantGfx{STFDEAD0}%
  \dtWantGfx{STFGOD0}%
  \dtQJ=0
  \dtWhile{\dtQJ<10}{%
    \edef\dt@next{\noexpand\dtWantGfx{STTNUM\the\dtQJ}}\dt@next
    \advance\dtQJ by1\relax }%
  \dtQJ=0
  \dtWhile{\dtQJ<5}{%
    \dtQK=0
    \dtWhile{\dtQK<3}{%
      \edef\dt@next{\noexpand\dtWantGfx{STFST\the\dtQJ\the\dtQK}}\dt@next
      \advance\dtQK by1\relax }%
    %% The four faces that only damage can reach: the two turns toward
    %% whatever hit you, the ouch, the evil grin at a new weapon and the
    %% rampage.  Five more per pain level, and they go through the same
    %% single directory pass -- a lookup each would be twenty more scans of
    %% a 2928-entry lump directory.
    \edef\dt@next{\noexpand\dtWantGfx{STFTL\the\dtQJ0}}\dt@next
    \edef\dt@next{\noexpand\dtWantGfx{STFTR\the\dtQJ0}}\dt@next
    \edef\dt@next{\noexpand\dtWantGfx{STFOUCH\the\dtQJ}}\dt@next
    \edef\dt@next{\noexpand\dtWantGfx{STFEVL\the\dtQJ}}\dt@next
    \edef\dt@next{\noexpand\dtWantGfx{STFKILL\the\dtQJ}}\dt@next
    \advance\dtQJ by1\relax }%
  %% The rest of the bar: the arms panel and its six numbers in two fonts,
  %% the three key icons, and the small digits the ammunition boxes are drawn
  %% from.  Twenty-two more pictures, and every one of them rides the same
  %% single pass over the directory -- a lookup each would be twenty-two
  %% more scans of 2928 entries.
  \dtWantGfx{STARMS}%
  \dtQJ=0
  \dtWhile{\dtQJ<3}{%
    \edef\dt@next{\noexpand\dtWantGfx{STKEYS\the\dtQJ}}\dt@next
    \advance\dtQJ by1\relax }%
  \dtQJ=0
  \dtWhile{\dtQJ<10}{%
    \edef\dt@next{\noexpand\dtWantGfx{STYSNUM\the\dtQJ}}\dt@next
    \advance\dtQJ by1\relax }%
  \dtQJ=2
  \dtWhile{\dtQJ<8}{%
    \edef\dt@next{\noexpand\dtWantGfx{STGNUM\the\dtQJ}}\dt@next
    \advance\dtQJ by1\relax }}

%%% ===============================================================  TEXTURES
%%% A Doom wall texture is not a picture: it is a recipe.  TEXTURE1 gives a
%%% size and a list of patches with offsets, PNAMES turns a patch number into
%%% a lump name, and the patch itself is the same column/post format
%%% \dt@readpic already reads for sprites.  Compositing once at load into
%%% dtX<t>@<u>@<v> gives the renderer a flat array to sample, so nothing in
%%% the inner loop has to know any of this.
%%%
%%% Only the textures this map's sidedefs actually name are built -- MAP01
%%% names 39 of the IWAD's 428 -- and their patches are resolved in ONE pass
%%% over the directory, for the same reason the sprite loader does it that
%%% way: \dtFindLump scans from a given entry and this WAD has 2928 of them.
\newcount\NTEX \NTEX=0
\newcount\dtXJ \newcount\dtXK \newcount\dtXN \newcount\dtXT
\newcount\dtXU \newcount\dtXV \newcount\dtXZ
\newcount\dtXW \newcount\dtXH \newcount\dtXA \newcount\dtXB
\newcount\dtXC \newcount\dtXD \newcount\dtXL
\newcount\dtXQ \newcount\dtXR \newcount\dtXOX \newcount\dtXOY
\newcount\dtXPI \newcount\dtXPC \newcount\dtXPW \newcount\dt@ub
\newcount\dtXNP \newcount\dtXMISS

%%% Lump and texture names are compared uppercased.  Doom does the same, and
%%% it is not a formality: MAP01's sidedefs name `doortrak' and `support2' in
%%% lower case while TEXTURE1 spells both in upper.  Matching raw bytes drops
%%% exactly those two textures and leaves holes in the map.
\def\dtUpHex{\gdef\dt@tnm{}\dt@uphA}
\def\dt@uphA#1#2{\ifx\relax#1\else
  \dt@ub=\numexpr\dtB{#1}{#2}\relax
  \ifnum\dt@ub>96 \ifnum\dt@ub<123 \advance\dt@ub by-32 \fi\fi
  \xdef\dt@tnm{\dt@tnm\hx{\the\dt@ub}}%
  \expandafter\dt@uphA\fi}
%% \dt@tnm holds a name; is it a real one?  A sidedef with no texture on a
%% surface spells it `-', and an unused slot is eight NULs.
\newif\ifdt@hastex
\def\dt@tfirst#1#2#3\relax{\def\dt@tf{#1#2}}
\def\dt@testname{\expandafter\dt@tfirst\dt@tnm\relax
  \dt@hastextrue
  \ifnum\pdfstrcmp{\dt@tf}{2D}=0 \dt@hastexfalse\fi
  \ifnum\pdfstrcmp{\dt@tf}{00}=0 \dt@hastexfalse\fi}

%%% Read the 8-byte name at offset #1 and normalise it into \dt@tnm.
\def\dtReadName#1{%
  \edef\dt@rnm{\pdffiledump offset #1 length 8 {\dtWADFILE}}%
  \expandafter\dtUpHex\dt@rnm\relax\relax}

%%% \dtWantTex : \dt@tnm names a texture this map uses.  Leaves its index in
%%% \dtXN, allocating one on first sight.  The index is handed out before the
%%% texture is known to exist; an unresolved one keeps width 0 and reads as
%%% absent, so a sidedef naming a texture the WAD does not have simply falls
%%% back to the procedural table rather than failing the load.
\def\dtWantTex{%
  \ifcsname dtXi\dt@tnm\endcsname\else
    \expandafter\xdef\csname dtXi\dt@tnm\endcsname{\the\NTEX}%
    \expandafter\xdef\csname dtXw\the\NTEX\endcsname{0}%
    \expandafter\xdef\csname dtXh\the\NTEX\endcsname{0}%
    \global\advance\NTEX by1
    \dt@swpair
  \fi
  \dtXN=\csname dtXi\dt@tnm\endcsname\relax}

%%% ---- switch faces --------------------------------------------------------
%%% Doom holds the SW1/SW2 pairing in the engine, not in the WAD -- vanilla
%%% has no SWITCHES lump -- but the pairing it holds is the naming convention
%%% itself, so it can be recovered from the name.  A texture whose name
%%% begins SW1 is wanted together with its SW2 twin (and the reverse), and
%%% the two indices are recorded in dtXs pointing at each other.
%%%
%%% Doing it here rather than when the switch is pressed is what makes the
%%% pressed face exist at all: only the names a sidedef mentions are ever
%%% composited, and a map naturally mentions only the unpressed one.
%%%
%%% The name is already hex, uppercased and NUL-padded to eight bytes, so
%%% "SW" is the four characters 5357 and the digit is the two after it; the
%%% counterpart differs in exactly one of the sixteen.  Comparison is
%%% character by character with \if rather than \ifnum, because a name whose
%%% hex contains A-F is not a number and \ifnum would raise an error on it
%%% (`JOINT2' would).  The ZZZZZZ sentinel makes the six-argument grab safe
%%% for a name shorter than three characters; Z is not a hex digit, so the
%%% tail can be recovered by delimiting on it.
\def\dt@swpair{%
  \begingroup
    \edef\dt@swa{\dt@tnm}%
    \expandafter\dt@swsplit\dt@tnm ZZZZZZ\relax
  \endgroup}
\def\dt@swtail#1ZZZZZZ{#1}
\def\dt@swsplit#1#2#3#4#5#6#7\relax{%
  \if#15\if#23\if#35\if#47
    \if#53
      \if#61 \dt@swmake{32}{#7}\fi
      \if#62 \dt@swmake{31}{#7}\fi
    \fi
  \fi\fi\fi\fi}
%%% #1 is the toggled digit as hex, #2 the rest of the name with the sentinel
%%% still on it.  \dtWantTex recurses back into \dt@swpair for the twin, but
%%% only on the branch that allocates, so the twin's own pairing finds this
%%% name already present and the recursion stops one level down.
\def\dt@swmake#1#2{%
  \edef\dt@tnm{5357#1\dt@swtail#2}%
  \dtWantTex
  \expandafter\xdef\csname dtXs\csname dtXi\dt@swa\endcsname\endcsname
    {\the\dtXN}%
  \expandafter\xdef\csname dtXs\the\dtXN\endcsname
    {\csname dtXi\dt@swa\endcsname}}
%% -1 when the name is `-' or empty, which is what \dtAddSide stores.
\def\dtTexRef#1{\dtReadName{#1}\dt@testname
  \ifdt@hastex \dtWantTex \else \dtXN=-1 \fi}
%% Want a texture by its literal name rather than by a name in the file.
%% \dtNameHex does not fold case, so the argument must already be upper.
\def\dtWantTexName#1{\dtNameHex{#1}\edef\dt@tnm{\dt@target}\dtWantTex}

%%% ---- PNAMES: a count, then one 8-byte patch name each --------------------
\def\dt@loadpnames{%
  \dtXNP=0
  \dtFindLump{PNAMES}{0}%
  \ifnum\dtFOUND<0
    \message{[wad] no PNAMES -- walls stay on the procedural tables.}%
  \else
    \dtXNP=\dtL{\the\dtLPOS}\relax
    \dtXJ=0
    \dtWhile{\dtXJ<\dtXNP}{%
      \dtReadName{\number\numexpr\dtLPOS+4+8*\dtXJ\relax}%
      \expandafter\xdef\csname dtXp\the\dtXJ\endcsname{\dt@tnm}%
      \advance\dtXJ by1\relax }%
  \fi}

%%% ---- TEXTURE1 / TEXTURE2: a count, an offset each, then maptexture_t -----
%%% maptexture_t is name[8], masked int32, width, height, columndirectory
%%% int32, patchcount -- 22 bytes -- then patchcount mappatch_t of 10.  Only
%%% the record's offset is kept here; the patch list is re-read at composite
%%% time rather than being copied into arrays it would only be read from once.
\def\dt@scantex#1{%
  \dtFindLump{#1}{0}%
  \ifnum\dtFOUND<0\else
    \dtXA=\dtLPOS
    \dtXC=\dtL{\the\dtXA}\relax
    \dtXJ=0
    \dtWhile{\dtXJ<\dtXC}{%
      \dtXB=\numexpr\dtXA+\dtL{\number\numexpr\dtXA+4+4*\dtXJ\relax}\relax
      \dtReadName{\the\dtXB}%
      \ifcsname dtXi\dt@tnm\endcsname
        \dtXT=\csname dtXi\dt@tnm\endcsname\relax
        \ifnum\csname dtXw\the\dtXT\endcsname=0
          \expandafter\xdef\csname dtXo\the\dtXT\endcsname{\the\dtXB}%
          \expandafter\xdef\csname dtXw\the\dtXT\endcsname
            {\dtS{\number\numexpr\dtXB+12\relax}}%
          \expandafter\xdef\csname dtXh\the\dtXT\endcsname
            {\dtS{\number\numexpr\dtXB+14\relax}}%
          %% v wraps at the texture's own height, so precompute the modulus
          %% the render loop subtracts rather than multiplying per column.
          \expandafter\xdef\csname dtXm\the\dtXT\endcsname
            {\number\numexpr\csname dtXh\the\dtXT\endcsname*65536\relax}%
        \fi
      \fi
      \advance\dtXJ by1\relax }%
  \fi}

%%% ---- which patch lumps are needed, and where they are --------------------
%%% Mark every patch named by a wanted texture, then walk the directory once.
\def\dt@wantpatches{%
  \dtXT=0
  \dtWhile{\dtXT<\NTEX}{%
    \ifnum\csname dtXw\the\dtXT\endcsname>0
      \dtXB=\csname dtXo\the\dtXT\endcsname\relax
      \dtXC=\dtUu{\number\numexpr\dtXB+20\relax}\relax
      \dtXJ=0
      \dtWhile{\dtXJ<\dtXC}{%
        \dtXPI=\dtUu{\number\numexpr\dtXB+22+10*\dtXJ+4\relax}\relax
        \ifnum\dtXPI<\dtXNP
          \expandafter\gdef\csname dtXq\csname dtXp\the\dtXPI\endcsname
            \endcsname{1}%
        \fi
        \advance\dtXJ by1\relax }%
    \fi
    \advance\dtXT by1\relax }}
\def\dt@findpatches{%
  \dtXJ=0
  \dtWhile{\dtXJ<\dtNUML}{%
    \dtReadName{\number\numexpr\dtDIR+16*\dtXJ+8\relax}%
    \ifcsname dtXq\dt@tnm\endcsname
      \ifcsname dtXl\dt@tnm\endcsname\else
        \expandafter\xdef\csname dtXl\dt@tnm\endcsname
          {\dtL{\number\numexpr\dtDIR+16*\dtXJ\relax}}%
      \fi
    \fi
    \advance\dtXJ by1\relax }}

%%% ---- compositing ---------------------------------------------------------
%%% A composited texel is stored EIGHT TO A CONTROL SEQUENCE, as
%%% dtX<t>@<g>@<v> holding the eight texels at u = 8g..8g+7 of row v.
%%%
%%% This is not a space nicety, it is the difference between the project
%%% fitting in pdfTeX and not.  One control sequence per texel is 420k
%%% strings for MAP01's 39 textures, against a pool of 497k under pdftex and
%%% 476k under pdflatex -- so the array alone was 92% of the budget and the
%%% first flat to be loaded killed the run with "TeX capacity exceeded
%%% [number of strings]".  Packed, the same textures cost 53k.
%%%
%%% Eight along u rather than along v because u is FIXED for a whole column
%%% span: \dt@usetex settles both the group and which of the eight to take
%%% once per column, so the inner loop pays one extra macro call and no
%%% arithmetic at all.  Packing along v would have to split the row per pixel.
%%%
%%% That in turn forces compositing to run COLUMN MAJOR, in groups of eight,
%%% because a group needs eight columns of one row at once.  Doom's order is
%%% patch major -- draw patch 1 everywhere, then patch 2 over it -- so this
%%% inverts it: for each column every patch that covers it writes into a
%%% scratch column, later patches overwriting earlier ones, which is the same
%%% semantics reached from the other side.
%%%
%%% The scratch dtZZ<k>@<v> is REDEFINED for every group rather than
%%% allocated per group.  Redefining a control sequence makes no new string,
%%% and that is the whole trick: 8 x (h+1) scratch names serve the entire
%%% texture set.
\newcount\dtXGG \newcount\dtXGK \newcount\dtXGU \newcount\dtXZE

%% Clearing is what makes a hole read as 00 rather than as the previous
%% group's pixel, so it replaces \dt@fillholes and its texel counter: those
%% only ran for the 2 textures of 39 that do not tile exactly, and there is
%% no equivalent of "did anything write here" once the array is packed.
\def\dt@zclear{%
  \dtXGK=0
  \dtWhile{\dtXGK<8}{%
    \dtXV=0
    \dtWhile{\dtXV<\dtXZE}{%
      \expandafter\gdef\csname dtZZ\the\dtXGK @\the\dtXV\endcsname{00}%
      \advance\dtXV by1\relax }%
    \advance\dtXGK by1\relax }}

%% One patch column's posts into scratch column \dtXGK.  A post that lies
%% wholly inside the texture -- almost all of them -- skips the bounds test
%% per pixel.  The guarded form still has to consume the clipped characters,
%% so it cannot simply stop early.
\def\dt@txrunF#1#2{%
  \expandafter\xdef\csname dtZZ\the\dtXGK @\the\dtXV\endcsname{#1#2}%
  \advance\dtXV by1
  \ifnum\dtXV<\dtXZ \expandafter\dt@txrunF\fi}
\def\dt@txrunG#1#2{%
  \ifnum\dtXV>-1 \ifnum\dtXV<\dtXH
    \expandafter\xdef\csname dtZZ\the\dtXGK @\the\dtXV\endcsname{#1#2}%
  \fi\fi
  \advance\dtXV by1
  \ifnum\dtXV<\dtXZ \expandafter\dt@txrunG\fi}
\def\dt@txposts{%
  \dtXD=\dtBy{\the\dtXR}\relax
  \ifnum\dtXD<255
    \advance\dtXPC by1
    \dtXL=\dtBy{\number\numexpr\dtXR+1\relax}\relax
    \ifnum\dtXL>0
      \dtXV=\numexpr\dtXOY+\dtXD\relax
      \dtXZ=\numexpr\dtXV+\dtXL\relax
      \edef\dt@trun{\pdffiledump offset \number\numexpr\dtXR+3\relax\space
        length \the\dtXL\space{\dtWADFILE}}%
      \ifnum\dtXV<0
        \expandafter\dt@txrunG\dt@trun
      \else
        \ifnum\dtXZ>\dtXH
          \expandafter\dt@txrunG\dt@trun
        \else
          \expandafter\dt@txrunF\dt@trun
        \fi
      \fi
    \fi
    \dtXR=\numexpr\dtXR+4+\dtXL\relax
    \ifnum\dtXPC<128 \dt@txposts \fi
  \fi}

%% Every patch that covers texture column \dtXGU, in file order.
\def\dt@zcolumn{%
  \dtXJ=0
  \dtWhile{\dtXJ<\dtXC}{%
    \dtXA=\numexpr\dtXB+22+10*\dtXJ\relax
    \dtXOX=\dtS{\the\dtXA}\relax
    \dtXOY=\dtS{\number\numexpr\dtXA+2\relax}\relax
    \dtXPI=\dtUu{\number\numexpr\dtXA+4\relax}\relax
    \ifnum\dtXPI<\dtXNP
      \edef\dt@pnm{\csname dtXp\the\dtXPI\endcsname}%
      \ifcsname dtXl\dt@pnm\endcsname
        \dtXQ=\csname dtXl\dt@pnm\endcsname\relax
        \dtXPW=\dtS{\the\dtXQ}\relax
        \dtXK=\numexpr\dtXGU-\dtXOX\relax
        \ifnum\dtXK>-1 \ifnum\dtXK<\dtXPW
          \dtXR=\numexpr\dtXQ+\dtL{%
            \number\numexpr\dtXQ+8+4*\dtXK\relax}\relax
          \dtXPC=0
          \dt@txposts
        \fi\fi
      \fi
    \fi
    \advance\dtXJ by1\relax }}

%% Row h repeats row 0, for the same reason the procedural table has 65 rows:
%% the span expansion reads v with \numexpr, whose division ROUNDS, so a v
%% just short of the wrap asks for row h.
\def\dt@zpack{%
  \dtXGK=0
  \dtWhile{\dtXGK<8}{%
    \expandafter\xdef\csname dtZZ\the\dtXGK @\the\dtXH\endcsname
      {\csname dtZZ\the\dtXGK @0\endcsname}%
    \advance\dtXGK by1\relax }%
  \dtXV=0
  \dtWhile{\dtXV<\dtXZE}{%
    \expandafter\xdef\csname dtX\the\dtXT @\the\dtXGG @\the\dtXV\endcsname{%
      {\csname dtZZ0@\the\dtXV\endcsname}{\csname dtZZ1@\the\dtXV\endcsname}%
      {\csname dtZZ2@\the\dtXV\endcsname}{\csname dtZZ3@\the\dtXV\endcsname}%
      {\csname dtZZ4@\the\dtXV\endcsname}{\csname dtZZ5@\the\dtXV\endcsname}%
      {\csname dtZZ6@\the\dtXV\endcsname}{\csname dtZZ7@\the\dtXV\endcsname}}%
    \advance\dtXV by1\relax }}

\def\dt@compose{%
  \dtXW=\csname dtXw\the\dtXT\endcsname\relax
  \dtXH=\csname dtXh\the\dtXT\endcsname\relax
  \dtXB=\csname dtXo\the\dtXT\endcsname\relax
  \dtXC=\dtUu{\number\numexpr\dtXB+20\relax}\relax
  \dtXZE=\numexpr\dtXH+1\relax
  \dtXGG=0
  \dtWhile{\numexpr\dtXGG*8\relax<\dtXW}{%
    \dt@zclear
    \dtXGK=0
    \dtWhile{\dtXGK<8}{%
      \dtXGU=\numexpr\dtXGG*8+\dtXGK\relax
      \ifnum\dtXGU<\dtXW \dt@zcolumn \fi
      \advance\dtXGK by1\relax }%
    \dt@zpack
    \advance\dtXGG by1\relax }}

%%% ---- masked middles (M19) ------------------------------------------------
%%% A middle texture on a TWO-SIDED line is the one surface in Doom that is
%%% drawn with its holes: the room beyond shows through the bars.  \dt@compose
%%% above cannot serve it, and not by oversight -- \dt@zclear fills every
%%% untouched texel with palette 0 precisely so that an undefined one cannot
%%% raise `Missing \endcsname' in the inner loop, and that is exactly the
%%% information transparency needs.
%%%
%%% So those textures keep their PATCH's post chain as well as the composite,
%%% under the picture key MK<t>, and the sprite pass draws them -- \dtSprRun
%%% already walks posts, already resamples source driven and already splices
%%% one band per post.  The feature needs no new inner loop at all.
%%%
%%% Which textures: marked here rather than at \dtTexRef, because "two-sided"
%%% is a property of the LINE and a sidedef does not know what it is attached
%%% to.  This runs after LINEDEFS, which is why \dtLoadTextures is where it
%%% is called from.
\newcount\dtXMK \newcount\dtXNM \dtXNM=0
\def\dt@markmasked{%
  \dtXMK=0
  \dtWhile{\dtXMK<\NLINE}{%
    \ifcsname dtLB\the\dtXMK\endcsname
      \ifnum\csname dtLB\the\dtXMK\endcsname<0 \else
      \ifnum\csname dtLF\the\dtXMK\endcsname<0 \else
        \dt@markmaskside{dtLR\the\dtXMK}%
        \dt@markmaskside{dtLL\the\dtXMK}%
      \fi\fi
    \fi
    \advance\dtXMK by1\relax }}
\def\dt@markmaskside#1{%
  \ifcsname #1\endcsname
    \ifnum\csname #1\endcsname<0 \else
      \edef\dt@mkt{\csname dtYm\csname #1\endcsname\endcsname}%
      \ifnum\dt@mkt<0 \else
        \expandafter\gdef\csname dtXk\dt@mkt\endcsname{1}%
      \fi
    \fi
  \fi}

%%% Two ways to get them.  A texture that is ONE patch at (0,0) -- which is
%%% every masked middle with a hole in it that the commercial maps use;
%%% MAP01's two are MIDBARS3 out of RW45_1 and BRNSMAL1 out of W111_2 -- has
%%% its posts already, in the patch, and \dt@readpic takes them as they are.
%%%
%%% Anything else is built (M23).  Doom itself gives up here:
%%% R_GenerateComposite keeps posts only for a column a single patch covers
%%% and draws the rest as garbage, which is the "medusa" every mapper learns
%%% to avoid.  There is no reason to copy that.  Each column goes through
%%% \dt@zcolumn exactly as the opaque composite did -- every patch that
%%% covers it, later ones over earlier -- but into a scratch column cleared
%%% to a marker instead of to palette 0, and the rows between markers are
%%% gathered back into posts.  Doom II names three such middles (PIPE4,
%%% SKSPINE2 and TEKBRON1, each two patches side by side), and until this
%%% they were not drawn at all.
%%%
%%% \dt@compose has just run for this texture, so \dtXW, \dtXH, \dtXB and
%%% \dtXC are already its own.  A patch narrower or shorter than the texture
%%% is fine either way: where nothing was drawn there is no post, which
%%% reads as transparent.
\newif\ifdt@mkone
\newcount\dtXMU \newcount\dtXMV \newcount\dtXMS
\def\dt@maskread{%
  \dt@mkonefalse
  \ifnum\dtXC=1
    \dtXA=\numexpr\dtXB+22\relax
    \dtXOX=\dtS{\the\dtXA}\relax
    \dtXOY=\dtS{\number\numexpr\dtXA+2\relax}\relax
    \ifnum\dtXOX=0 \ifnum\dtXOY=0
      \dtXPI=\dtUu{\number\numexpr\dtXA+4\relax}\relax
      \ifnum\dtXPI<\dtXNP
        \edef\dt@pnm{\csname dtXp\the\dtXPI\endcsname}%
        \ifcsname dtXl\dt@pnm\endcsname
          \dtLPOS=\csname dtXl\dt@pnm\endcsname\relax
          \edef\dt@next{\noexpand\dt@readpic{MK\the\dtXT}{0}}\dt@next
          \dt@mkonetrue
        \fi
      \fi
    \fi\fi
  \fi
  \ifdt@mkone\else \dt@maskbuild \fi
  \expandafter\gdef\csname dtXn\the\dtXT\endcsname{1}%
  \global\advance\dtXNM by1 }

%% The marker is two characters that are not hex, so no pixel can be
%% mistaken for it; \ifx compares the two macros' bodies and both were typed
%% in this file, so the catcodes agree.  \dtXV belongs to \dt@txposts, which
%% \dt@zcolumn runs, so the gathering walk has a counter of its own.
\def\dt@mkhole{--}
\def\dt@maskbuild{%
  \expandafter\xdef\csname dtPwMK\the\dtXT\endcsname{\the\dtXW}%
  \expandafter\xdef\csname dtPhMK\the\dtXT\endcsname{\the\dtXH}%
  \expandafter\gdef\csname dtPlMK\the\dtXT\endcsname{0}%
  \expandafter\gdef\csname dtPtMK\the\dtXT\endcsname{0}%
  \dtXMU=0
  \dtWhile{\dtXMU<\dtXW}{%
    \dtXV=0
    \dtWhile{\dtXV<\dtXH}{%
      \expandafter\gdef\csname dtZZ0@\the\dtXV\endcsname{--}%
      \advance\dtXV by1\relax }%
    \dtXGK=0 \dtXGU=\dtXMU
    \dt@zcolumn
    \gdef\dt@post{}\gdef\dt@mrun{}%
    \dtXMS=-1 \dtXMV=0
    \dtWhile{\dtXMV<\dtXH}{%
      \edef\dt@mpx{\csname dtZZ0@\the\dtXMV\endcsname}%
      \ifx\dt@mpx\dt@mkhole
        \dt@maskflush
      \else
        \ifnum\dtXMS<0 \dtXMS=\dtXMV \fi
        \xdef\dt@mrun{\dt@mrun\dt@mpx}%
      \fi
      \advance\dtXMV by1\relax }%
    \dt@maskflush
    \expandafter\xdef\csname dtQMK\the\dtXT @\the\dtXMU\endcsname
      {\dt@post{-1}{0}{}}%
    \advance\dtXMU by1\relax }}
%% One run of drawn rows becomes one post, in \dt@readpic's own shape.
\def\dt@maskflush{%
  \ifnum\dtXMS<0 \else
    \xdef\dt@post{\dt@post{\the\dtXMS}{\number\numexpr\dtXMV-\dtXMS\relax}%
      {\dt@mrun}}%
    \gdef\dt@mrun{}\dtXMS=-1
  \fi}

%%% Diagnostic: put a texture that is already loaded through \dt@maskbuild,
%%% whichever way its posts were first got.  The commercial maps name no
%%% texture that is both several patches AND full of holes, so the only way
%%% to watch the builder find a hole is to hand it MIDBARS3 -- whose posts
%%% the patch has already supplied -- and require the same picture back.
\def\dtMaskRebuild#1{%
  \dtNameHex{#1}%
  \ifcsname dtXi\dt@target\endcsname
    \dtXT=\csname dtXi\dt@target\endcsname\relax
    \ifnum\csname dtXw\the\dtXT\endcsname>0
      \dtXW=\csname dtXw\the\dtXT\endcsname\relax
      \dtXH=\csname dtXh\the\dtXT\endcsname\relax
      \dtXB=\csname dtXo\the\dtXT\endcsname\relax
      \dtXC=\dtUu{\number\numexpr\dtXB+20\relax}\relax
      \dt@maskbuild
      \message{[maskrebuild] #1 \the\dtXW x\the\dtXH}%
    \fi
  \fi}

%%% Diagnostic: one column of a masked middle's posts, so a harness can build
%%% the same column from the WAD and compare -- \dtTexRow's argument, for the
%%% one array \dtTexRow cannot see.
\def\dtMaskCol#1#2{%
  \dtNameHex{#1}%
  \ifcsname dtXi\dt@target\endcsname
    \dtXT=\csname dtXi\dt@target\endcsname\relax
    \ifcsname dtQMK\the\dtXT @#2\endcsname
      \message{[maskcol] #1 #2 \csname dtQMK\the\dtXT @#2\endcsname}%
    \else
      \message{[maskcol] #1 #2 opaque}%
    \fi
  \else
    \message{[maskcol] #1 #2 missing}%
  \fi}

%%% Diagnostic: print one row of a composited texture as hex bytes, so that a
%%% harness can composite the same texture by its own route and compare.  It
%%% is the only way to check the composite itself rather than its effect --
%%% a wrong PNAMES index, a patch placed at the wrong origin, a post decoded
%%% one byte out or a hole filled that should have been painted all produce a
%%% picture that renders perfectly well and is simply not the right one.
%%%
%%% One \message, not one per column: pdfTeX hard-wraps terminal output at
%%% max_print_line without inserting or dropping characters, so a single
%%% string survives deleting the newlines while separate messages would have
%%% spaces between them.
\def\dtTexRow#1#2{%
  \dtNameHex{#1}%
  \ifcsname dtXi\dt@target\endcsname
    \dtXT=\csname dtXi\dt@target\endcsname\relax
    \dtXW=\csname dtXw\the\dtXT\endcsname\relax
    \dtXH=\csname dtXh\the\dtXT\endcsname\relax
    %% Sampled one texel at a time THROUGH \dt@pick, the same way \WallTex
    %% reaches a texel, rather than by taking whole groups.  Taking groups
    %% would be shorter and would leave the pick -- the half of the packing
    %% that runs per column -- with nothing checking it at all: a \dt@usetex
    %% that always chose the first of the eight would draw every wall as
    %% eight-fold vertical stripes and pass every check in this file.
    \gdef\dt@row{}%
    \dtXU=0
    \dtWhile{\dtXU<\dtXW}{%
      %% Through \dt@usetex, not by splitting u here.  Working out the group
      %% and the pick independently would check the STORAGE and leave the
      %% half of the packing that runs per column unguarded: a \dt@usetex
      %% that always chose the first texel of the eight draws every wall as
      %% eight-fold vertical stripes and passes every render check in
      %% check.py.  Confirmed by making it do exactly that.
      %% \rd is the u \dt@usetex wraps, \vstepr the step \dt@vsetup reduces.
      \rd=\dtXU \vstepr=65536
      \dt@usetex{\the\dtXT}%
      \xdef\dt@row{\dt@row
        \expandafter\expandafter\expandafter\dt@pick
          \csname\dt@ux#2\endcsname}%
      \advance\dtXU by1\relax }%
    \message{[texrow] #1 #2 \the\dtXW\space\the\dtXH\space\dt@row}%
  \else
    \message{[texrow] #1 #2 missing}%
  \fi}

%%% ==================================================================  FLATS
%%% A flat is the one Doom graphic that is not a picture and not a recipe:
%%% 4096 raw bytes, 64 rows of 64, no posts, no offsets, no PNAMES.  So the
%%% loader is a fraction of \dtLoadTextures -- the only thing it borrows is
%%% the shape, resolving every wanted name in ONE pass over the directory
%%% between the F_START and F_END markers rather than a \dtFindLump each.
%%%
%%% Stored as dtFL<t>@<v>@<g>, a row v split into eight groups of eight
%%% texels, indexed [row][column] because that is how the span arithmetic
%%% arrives: Doom's R_DrawSpan reads
%%% source[((yfrac>>10)&4032) + ((xfrac>>16)&63)], i.e. world y picks the row
%%% and world x the column.  Getting that pair the wrong way round transposes
%%% every floor in the map, which on a symmetric flat like FLOOR4_8 is
%%% invisible -- see check.py's asymmetric-flat probe.
%%%
%%% Eight to a control sequence rather than one, because THE STRING POOL IS
%%% THE BINDING RESOURCE HERE.  39 composited wall textures already hold
%%% ~457k of pdfTeX's 497k strings; a texel each would want 127k more and the
%%% load dies with "TeX capacity exceeded [number of strings]" halfway
%%% through the second flat.  Packing costs an eight-argument grab per pixel
%%% and spends `words of memory', which this build has 2.8M of going spare.
%%% Anything that adds a per-texel array from here on has the same problem
%%% and the same answer.
\newcount\NFLAT \NFLAT=0
\newcount\dtFN \newcount\dtFT \newcount\dtFU \newcount\dtFV
\newcount\dtFJ \newcount\dtFB \newcount\dtFE \newcount\dtFP
\newcount\dtFMISS \newcount\dtHASSKY \dtHASSKY=0

%%% F_SKY1 is not a flat.  It is the marker Doom uses to say "this plane is
%%% the sky", and its 4096 bytes are never drawn -- drawing them gives a
%%% ceiling of red rock that stretches to infinity at the horizon.  Ten of
%%% MAP01's 59 sectors name it, so this is not an edge case.
%% \dtNameHex walks its argument a character at a time and does not expand
%% it, so the name has to be written out here rather than kept in a macro.
\dtNameHex{F_SKY1}\edef\dt@skyhex{\dt@target}

%%% \dtWantFlat : \dt@tnm names a flat this map uses.  Same contract as
%%% \dtWantTex -- the index is handed out before the lump is known to exist,
%%% and dtFw stays 0 until the F_START pass actually reads one, so a sector
%%% naming a flat the WAD has not got falls back to its procedural ramp
%%% instead of failing the load.
\def\dtWantFlat{%
  \ifcsname dtFi\dt@tnm\endcsname\else
    \expandafter\xdef\csname dtFi\dt@tnm\endcsname{\the\NFLAT}%
    \expandafter\xdef\csname dtFw\the\NFLAT\endcsname{0}%
    \global\advance\NFLAT by1
  \fi
  \dtFN=\csname dtFi\dt@tnm\endcsname\relax}

%% -1 when the name is absent, -2 when it is the sky.
\def\dtFlatRef#1{\dtReadName{#1}\dt@testname
  \ifdt@hastex
    \ifnum\pdfstrcmp{\dt@tnm}{\dt@skyhex}=0
      \dtFN=-2 \global\dtHASSKY=1
    \else \dtWantFlat \fi
  \else \dtFN=-1 \fi}

%%% One row of 64 bytes is ONE \pdffiledump, split into eight groups of eight
%%% here.  The per-pixel alternative would be 4096 file dumps a flat.  Two
%%% stages because a macro takes nine arguments and a group is sixteen
%%% characters; between them they consume the row exactly.
\def\dt@flrunA#1#2#3#4#5#6#7#8{%
  \def\dt@fgrp{{#1#2}{#3#4}{#5#6}{#7#8}}\dt@flrunB}
\def\dt@flrunB#1#2#3#4#5#6#7#8{%
  \expandafter\xdef\csname dtFL\the\dtFT @\the\dtFV @\the\dtFU\endcsname
    {\dt@fgrp{#1#2}{#3#4}{#5#6}{#7#8}}%
  \advance\dtFU by1
  \ifnum\dtFU<8 \expandafter\dt@flrunA\fi}
\def\dt@readflat{%
  \dtFV=0
  \dtWhile{\dtFV<64}{%
    \edef\dt@frow{\pdffiledump offset
      \number\numexpr\dtFP+64*\dtFV\relax\space length 64 {\dtWADFILE}}%
    \dtFU=0
    \expandafter\dt@flrunA\dt@frow
    \advance\dtFV by1\relax }}

%%% Diagnostic: print one row of a loaded flat as hex bytes, so a harness can
%%% read the same 64 bytes out of the WAD itself and compare.  Same reasoning
%%% as \dtTexRow -- a flat that is transposed, off by a row, or read from the
%%% wrong lump renders perfectly plausibly and is simply not the right one.
%%% One \message, because pdfTeX hard-wraps terminal output mid-token.
\def\dtFlatRow#1#2{%
  \dtNameHex{#1}%
  \ifcsname dtFi\dt@target\endcsname
    \dtFT=\csname dtFi\dt@target\endcsname\relax
    \gdef\dt@frw{}%
    \dtFU=0
    \dtWhile{\dtFU<8}{%
      \xdef\dt@frw{\dt@frw\csname dtFL\the\dtFT @#2@\the\dtFU\endcsname}%
      \advance\dtFU by1\relax }%
    \message{[flatrow] #1 #2 \dt@frw}%
  \else
    \message{[flatrow] #1 #2 missing}%
  \fi}

\def\dtLoadFlats{%
  %% \DTNOTEX means "keep every surface on the procedural shade tables", and
  %% the palette-ramp invariant in INTERNALS section 6 is stated about exactly
  %% that render -- so it has to cover planes too, or the check goes red for
  %% the right reason and the wrong surface.  \DTNOFLAT turns off only the
  %% planes, which is what makes a flats-on/flats-off diff possible.
  \expandafter\ifx\csname DTNOTEX\endcsname\relax\else \global\NFLAT=0 \fi
  \expandafter\ifx\csname DTNOFLAT\endcsname\relax\else
    \global\NFLAT=0
    \message{[wad] \string\DTNOFLAT: planes stay on the procedural ramps}%
  \fi
  \ifnum\NFLAT>0
    \dtFindLump{F_START}{0}%
    \ifnum\dtFOUND<0
      \message{[wad] no F_START -- planes stay on the procedural ramps.}%
      \global\NFLAT=0
    \else
      \dtFB=\dtFOUND
      \dtFE=\dtNUML
      \dtFindLump{F_END}{\the\dtFB}%
      \ifnum\dtFOUND<0\else \dtFE=\dtFOUND \fi
      \dtFJ=\numexpr\dtFB+1\relax
      \dtWhile{\dtFJ<\dtFE}{%
        \dtReadName{\number\numexpr\dtDIR+16*\dtFJ+8\relax}%
        \ifcsname dtFi\dt@tnm\endcsname
          \ifnum\csname dtFw\csname dtFi\dt@tnm\endcsname\endcsname=0
            \dtFT=\csname dtFi\dt@tnm\endcsname\relax
            \dtFP=\dtL{\number\numexpr\dtDIR+16*\dtFJ\relax}\relax
            \ifnum\dtL{\number\numexpr\dtDIR+16*\dtFJ+4\relax}<4096 \else
              \dt@readflat
              \expandafter\xdef\csname dtFw\the\dtFT\endcsname{1}%
            \fi
          \fi
        \fi
        \advance\dtFJ by1\relax }%
      \dtFMISS=0
      \dtFT=0
      \dtWhile{\dtFT<\NFLAT}{%
        \ifnum\csname dtFw\the\dtFT\endcsname=0 \advance\dtFMISS by1 \fi
        \advance\dtFT by1\relax }%
      \message{[wad] flats: \number\numexpr\NFLAT-\dtFMISS\relax\space read
        of \the\NFLAT\space named%
        \ifnum\dtFMISS>0 , \the\dtFMISS\space not between F_START and F_END\fi
        \ifnum\dtHASSKY>0 ; sky planes present\fi}%
    \fi
  \fi}

\def\dtLoadTextures{%
  %% \DTNOTEX keeps the walls on the procedural shade tables even with a real
  %% IWAD loaded.  check.py needs it: two of its oldest invariants are about
  %% how *procedurally shaded* surfaces behave, and neither can be stated
  %% about a wall carrying whatever colours the artist used.
  \expandafter\ifx\csname DTNOTEX\endcsname\relax\else
    \global\NTEX=0
    \message{[wad] \string\DTNOTEX: walls stay on the procedural tables}%
  \fi
  %% The sky is a wall texture, not a flat: Doom draws it as a vertical strip
  %% indexed by view angle, so it has to join the want list before TEXTURE1 is
  %% scanned.  Which one it is depends on the map -- Doom II uses SKY1 for
  %% MAP01-11, SKY2 for 12-20, SKY3 above that -- so \DTSKY overrides it.
  \expandafter\ifx\csname DTNOFLAT\endcsname\relax\else \global\dtHASSKY=0 \fi
  %% It composites like any other texture -- which it can only afford to do
  %% because dtX is packed: a 256x128 sky is 33k texels, and unpacked that is
  %% more than the whole string pool had spare.
  \ifnum\dtHASSKY>0 \ifnum\NTEX>0
    \expandafter\ifx\csname DTSKY\endcsname\relax\def\DTSKY{SKY1}\fi
    \expandafter\dtWantTexName\expandafter{\DTSKY}%
    \global\dtSKYT=\dtXN
  \fi\fi
  \ifnum\NTEX>0
    \dt@loadpnames
    \ifnum\dtXNP>0
      \dt@scantex{TEXTURE1}%
      \dt@scantex{TEXTURE2}%
      \dt@wantpatches
      \dt@findpatches
      %% Before the compose loop, because \dt@maskread runs inside it and
      %% needs to know which textures were asked for as masked middles.
      %% \DTNOMASK keeps every middle opaque, which is the world M18 left.
      \expandafter\ifx\csname DTNOMASK\endcsname\relax \dt@markmasked \fi
      \dtXMISS=0
      \dtXT=0
      \dtWhile{\dtXT<\NTEX}{%
        \ifnum\csname dtXw\the\dtXT\endcsname>0
          \dt@compose
          \ifcsname dtXk\the\dtXT\endcsname \dt@maskread \fi
        \else
          \advance\dtXMISS by1
        \fi
        \advance\dtXT by1\relax }%
      \message{[wad] textures: \number\numexpr\NTEX-\dtXMISS\relax\space
        composited of \the\NTEX\space named%
        \ifnum\dtXMISS>0 , \the\dtXMISS\space not in TEXTURE1\fi
        \ifnum\dtXNM>0 ; \the\dtXNM\space kept as masked middles\fi}%
    \fi
  \fi}

%%% ------------------------------------------------------------- main loader
\newcount\dtNV \newcount\dtNSD \newcount\dtNSE \newcount\dtNLD
\newcount\dtNSG \newcount\dtNSS \newcount\dtNND
\newcount\dtVP \newcount\dtSDP \newcount\dtSEP \newcount\dtLDP
\newcount\dtSGP \newcount\dtSSP \newcount\dtNODP
\newcount\dtMAP \newcount\dtJ \newcount\dtFRONT \newcount\dtBACK
\newcount\dtSC \newcount\dtRW \newcount\dtRF \newcount\dtRC
\newcount\dtTT \newcount\dtTF \newcount\dtNTH \newcount\dtTB
\newcount\dtLSP \newcount\dtLFL \newcount\dtPEG \newcount\dtPGB
\newcount\dtLTG \newcount\dtSCG

%%% Doom sectors carry texture names, not palette ramps, so pick a ramp per
%%% sector to keep adjacent rooms visually distinct.
%%%
%%% Not every palette row can serve as a ramp base.  The built-in palette is
%%% sixteen monotonic ramps of sixteen by construction, but a real PLAYPAL is
%%% not: rows 9, 10 and 14 each hold *two* eight-entry ramps, and rows 0 and 15
%%% mix unrelated colours.  The renderer emits shades 0..13 from a base, so
%%% basing a surface on one of those rows walks off the end of the ramp and the
%%% surface jumps in brightness halfway down.  These eleven bases -- 16, 32,
%%% 48, 64, 80, 96, 112, 128, 176, 192, 208 -- are the rows that descend
%%% monotonically across the whole range the renderer can emit.  Each table
%%% below is a permutation of exactly those eleven.
%%%
%%% The first three entries are load bearing: they reproduce the built-in
%%% map's own sector ramps, which is what lets check.py compare a parsed WAD
%%% against the hand-written level.  Reordering them does not "only" break
%%% pixel identity -- \dtEdges keys on the ramp index, so it breaks the
%%% geometry comparison too.  Change these three and you lose the strongest
%%% check in the project.
\def\dtRampW#1{\ifcase#1 64\or128\or192\or32\or176\or208\or16\or48\or80\or
  96\or112\else64\fi}
\def\dtRampF#1{\ifcase#1 32\or112\or80\or16\or208\or64\or176\or192\or48\or
  96\or128\else32\fi}
\def\dtRampC#1{\ifcase#1 176\or192\or96\or48\or16\or112\or208\or32\or64\or
  80\or128\else176\fi}
\def\dtSectorRamps#1{%
  \dtSC=#1\relax
  \tmpb=\dtSC \divide\tmpb by11 \multiply\tmpb by11
  \advance\dtSC by-\tmpb
  \dtRW=\dtRampW{\the\dtSC}\relax
  \dtRF=\dtRampF{\the\dtSC}\relax
  \dtRC=\dtRampC{\the\dtSC}\relax}

\def\dtLoadWAD#1#2{%
  \gdef\dtWADFILE{#1}%
  \ifnum\pdffilesize{#1}=0
    \message{^^J[wad] cannot open #1 -- keeping the built-in map.^^J}%
  \else
    \message{^^J[wad] #1 (\pdffilesize{#1} bytes)}%
    \dtNUML=\dtL{4}\relax
    \dtDIR=\dtL{8}\relax
    \message{[wad] \the\dtNUML\space lumps, directory at \the\dtDIR}%
    \dtFindLump{#2}{0}%
    \ifnum\dtFOUND<0
      \message{[wad] map #2 not found -- keeping the built-in map.}%
    \else
      \dtMAP=\dtFOUND
      %% The map's lumps follow its marker in a fixed order.
      \dtFindLump{VERTEXES}{\the\dtMAP}\dtVP=\dtLPOS \dtNV=\dtLSIZE
        \divide\dtNV by4
      \dtFindLump{SIDEDEFS}{\the\dtMAP}\dtSDP=\dtLPOS \dtNSD=\dtLSIZE
        \divide\dtNSD by30
      \dtFindLump{SECTORS}{\the\dtMAP}\dtSEP=\dtLPOS \dtNSE=\dtLSIZE
        \divide\dtNSE by26
      \dtFindLump{LINEDEFS}{\the\dtMAP}\dtLDP=\dtLPOS \dtNLD=\dtLSIZE
        \divide\dtNLD by14
      \dtFindLump{SEGS}{\the\dtMAP}\dtSGP=\dtLPOS \dtNSG=\dtLSIZE
        \divide\dtNSG by12
      \dtFindLump{SSECTORS}{\the\dtMAP}\dtSSP=\dtLPOS \dtNSS=\dtLSIZE
        \divide\dtNSS by4
      \dtFindLump{NODES}{\the\dtMAP}\dtNODP=\dtLPOS \dtNND=\dtLSIZE
        \divide\dtNND by28
      \message{[wad] #2: \the\dtNV\space verts, \the\dtNLD\space lines,
        \the\dtNSD\space sides, \the\dtNSE\space sectors, \the\dtNSG\space
        segs, \the\dtNSS\space subsectors, \the\dtNND\space nodes}%
      \ifnum\dtNND<1
        \message{[wad] no BSP in this map -- cannot render it.}%
      \else
        %% ---- VERTEXES: two int16
        \dtJ=0
        \dtWhile{\dtJ<\dtNV}{%
          \edef\dt@x{\dtS{\number\numexpr\dtVP+4*\dtJ\relax}}%
          \edef\dt@y{\dtS{\number\numexpr\dtVP+4*\dtJ+2\relax}}%
          \expandafter\expandafter\expandafter\dtAddVert
            \expandafter\expandafter\expandafter{\expandafter\dt@x
            \expandafter}\expandafter{\dt@y}%
          \advance\dtJ by1\relax }%
        %% ---- SIDEDEFS: x and y offset, three texture names, sector.
        %% Naming a texture here is what puts it on the want list; the
        %% composite is built later, once the whole list is known.
        \dtJ=0
        \dtWhile{\dtJ<\dtNSD}{%
          \dtTexRef{\number\numexpr\dtSDP+30*\dtJ+4\relax}%
          \edef\dt@sdu{\the\dtXN}%
          \dtTexRef{\number\numexpr\dtSDP+30*\dtJ+12\relax}%
          \edef\dt@sdl{\the\dtXN}%
          \dtTexRef{\number\numexpr\dtSDP+30*\dtJ+20\relax}%
          \edef\dt@sdm{\the\dtXN}%
          \edef\dt@next{\noexpand\dtAddSide
            {\dtUu{\number\numexpr\dtSDP+30*\dtJ+28\relax}}%
            {\dt@sdu}{\dt@sdl}{\dt@sdm}%
            {\dtS{\number\numexpr\dtSDP+30*\dtJ\relax}}%
            {\dtS{\number\numexpr\dtSDP+30*\dtJ+2\relax}}}%
          \dt@next
          \advance\dtJ by1\relax }%
        %% ---- SECTORS: floor, ceiling, light at byte 20
        \dtJ=0
        \dtWhile{\dtJ<\dtNSE}{%
          \edef\dt@f{\dtS{\number\numexpr\dtSEP+26*\dtJ\relax}}%
          \edef\dt@c{\dtS{\number\numexpr\dtSEP+26*\dtJ+2\relax}}%
          \edef\dt@l{\dtUu{\number\numexpr\dtSEP+26*\dtJ+20\relax}}%
          \dtSectorRamps{\the\dtJ}%
          \edef\dt@next{\noexpand\dtAddSector{\dt@f}{\dt@c}{\dt@l}%
            {\the\dtRW}{\the\dtRF}{\the\dtRC}}%
          \dt@next
          %% floorpic at byte 4, ceilingpic at byte 12.  Naming one here is
          %% what puts it on the want list; the 4096 bytes are read later,
          %% once the whole list is known.  The ramps above stay: a flat that
          %% does not resolve falls back to them.
          \dtFlatRef{\number\numexpr\dtSEP+26*\dtJ+4\relax}%
          \expandafter\xdef\csname dtAf\the\dtJ\endcsname{\the\dtFN}%
          \dtFlatRef{\number\numexpr\dtSEP+26*\dtJ+12\relax}%
          \expandafter\xdef\csname dtAc\the\dtJ\endcsname{\the\dtFN}%
          %% byte 24 is the sector tag: which tagged linedefs act on it.
          %% Only nonzero tags are stored; \dtSecTag reports 0 for the rest.
          \dtSCG=\dtS{\number\numexpr\dtSEP+26*\dtJ+24\relax}\relax
          \ifnum\dtSCG=0\else
            \edef\dt@next{\noexpand\dtAddSectorTag{\the\dtJ}{\the\dtSCG}}%
            \dt@next
          \fi
          %% byte 22 is the sector's own special: a light that strobes, a
          %% floor that burns, a secret.  Nonzero ones only, as with the tag.
          \dtSCG=\dtS{\number\numexpr\dtSEP+26*\dtJ+22\relax}\relax
          \ifnum\dtSCG=0\else
            \edef\dt@next{\noexpand\dtAddSectorSpec{\the\dtJ}{\the\dtSCG}}%
            \dt@next
          \fi
          \advance\dtJ by1\relax }%
        %% ---- LINEDEFS: front and back sector via the sidedefs
        \dtJ=0
        \dtWhile{\dtJ<\dtNLD}{%
          \dtFRONT=\dtUu{\number\numexpr\dtLDP+14*\dtJ+10\relax}\relax
          \dtBACK=\dtUu{\number\numexpr\dtLDP+14*\dtJ+12\relax}\relax
          \edef\dt@fs{\ifnum\dtFRONT<\dtNSD
            \csname dtYs\the\dtFRONT\endcsname\else-1\fi}%
          \edef\dt@bs{\ifnum\dtBACK=65535 -1\else
            \ifnum\dtBACK<\dtNSD \csname dtYs\the\dtBACK\endcsname
            \else-1\fi\fi}%
          \edef\dt@next{\noexpand\dtAddLine{\dt@fs}{\dt@bs}}%
          \dt@next
          %% The sidedefs themselves, and the two pegging bits of the flags
          %% word at byte 4.  Pegging decides where a texture's v origin
          %% sits, and getting it wrong is most visible on doors: an upper
          %% that is NOT unpegged tracks the moving ceiling, which is what
          %% makes the slab travel with the door instead of standing still.
          \dtLFL=\dtUu{\number\numexpr\dtLDP+14*\dtJ+4\relax}\relax
          \dtPEG=\dtLFL \divide\dtPEG by8
          \tmpd=\dtPEG \divide\tmpd by2 \multiply\tmpd by2
          \advance\dtPEG by-\tmpd
          \dtPGB=\dtLFL \divide\dtPGB by16
          \tmpd=\dtPGB \divide\tmpd by2 \multiply\tmpd by2
          \advance\dtPGB by-\tmpd
          \multiply\dtPGB by2 \advance\dtPEG by\dtPGB
          %% ML_SOUNDBLOCK, bit 6 of the same flags word: the flag a mapper
          %% sets to keep one room from hearing the next.  M20's noise flood
          %% crosses at most one of these, and only nonzero ones are stored,
          %% so a map without any costs nothing and the built-in level needs
          %% no changes.
          \tmpd=\dtLFL \divide\tmpd by64
          \dtPGB=\tmpd \divide\dtPGB by2 \multiply\dtPGB by2
          \advance\tmpd by-\dtPGB
          \ifnum\tmpd=1
            \edef\dt@next{\noexpand\dtAddLineSound{\the\dtJ}}\dt@next
          \fi
          \edef\dt@next{\noexpand\dtAddLineSide{\the\dtJ}%
            {\ifnum\dtFRONT<\dtNSD \the\dtFRONT\else-1\fi}%
            {\ifnum\dtBACK<\dtNSD \the\dtBACK\else-1\fi}%
            {\the\dtPEG}}%
          \dt@next
          %% byte 6 is the linedef special: what the line does when used or
          %% crossed.  Only nonzero ones are stored; \dtLineSpec reports 0 for
          %% the rest, so the built-in map needs no changes.
          \dtLSP=\dtUu{\number\numexpr\dtLDP+14*\dtJ+6\relax}\relax
          \ifnum\dtLSP>0
            \edef\dt@next{\noexpand\dtAddLineSpec{\the\dtJ}{\the\dtLSP}}%
            \dt@next
          \fi
          %% byte 8 is the tag.  A manual door acts on its own back sector and
          %% carries tag 0; every other special acts on the sectors carrying
          %% this number, however far away they are, so the tag is the only
          %% link between the two and cannot be recovered from the geometry.
          \dtLTG=\dtS{\number\numexpr\dtLDP+14*\dtJ+8\relax}\relax
          \ifnum\dtLTG=0\else
            \edef\dt@next{\noexpand\dtAddLineTag{\the\dtJ}{\the\dtLTG}}%
            \dt@next
          \fi
          \advance\dtJ by1\relax }%
        %% ---- SEGS: v1, v2, linedef, side
        \dtJ=0
        \dtWhile{\dtJ<\dtNSG}{%
          \edef\dt@a{\dtUu{\number\numexpr\dtSGP+12*\dtJ\relax}}%
          \edef\dt@b{\dtUu{\number\numexpr\dtSGP+12*\dtJ+2\relax}}%
          \edef\dt@ld{\dtUu{\number\numexpr\dtSGP+12*\dtJ+6\relax}}%
          \edef\dt@sd{\dtUu{\number\numexpr\dtSGP+12*\dtJ+8\relax}}%
          \edef\dt@next{\noexpand\dtAddSeg{\dt@a}{\dt@b}{\dt@ld}{\dt@sd}}%
          \dt@next
          %% byte 10: how far along the linedef this seg starts.  Without it
          %% every seg restarts its texture at u=0 and a wall split by the
          %% BSP shows a seam at each split.
          \edef\dt@next{\noexpand\dtAddSegOff{\the\dtJ}%
            {\dtS{\number\numexpr\dtSGP+12*\dtJ+10\relax}}}%
          \dt@next
          \advance\dtJ by1\relax }%
        %% ---- SSECTORS: seg count, first seg
        \dtJ=0
        \dtWhile{\dtJ<\dtNSS}{%
          \edef\dt@n{\dtUu{\number\numexpr\dtSSP+4*\dtJ\relax}}%
          \edef\dt@f{\dtUu{\number\numexpr\dtSSP+4*\dtJ+2\relax}}%
          \edef\dt@next{\noexpand\dtAddSSec{\dt@n}{\dt@f}}%
          \dt@next
          \advance\dtJ by1\relax }%
        %% ---- NODES: partition line and the two children.  The bounding
        %% ---- boxes at bytes 8..23 are skipped: this renderer prunes with
        %% ---- the open-column count instead of CheckBBox.
        \dtJ=0
        \dtWhile{\dtJ<\dtNND}{%
          \edef\dt@x{\dtS{\number\numexpr\dtNODP+28*\dtJ\relax}}%
          \edef\dt@y{\dtS{\number\numexpr\dtNODP+28*\dtJ+2\relax}}%
          \edef\dt@u{\dtS{\number\numexpr\dtNODP+28*\dtJ+4\relax}}%
          \edef\dt@v{\dtS{\number\numexpr\dtNODP+28*\dtJ+6\relax}}%
          \edef\dt@cr{\dtUu{\number\numexpr\dtNODP+28*\dtJ+24\relax}}%
          \edef\dt@cl{\dtUu{\number\numexpr\dtNODP+28*\dtJ+26\relax}}%
          \edef\dt@next{\noexpand\dtAddNode{\dt@x}{\dt@y}{\dt@u}{\dt@v}%
            {\dt@cr}{\dt@cl}}%
          \dt@next
          \advance\dtJ by1\relax }%
        %% Doom's root node is the last one built.
        \DTROOT=\numexpr\NNODE-1\relax
        \message{[wad] loaded: \the\NVERT\space verts, \the\NLINE\space lines,
          \the\NSEC\space sectors, \the\NSEG\space segs, \the\NSSEC\space
          subsectors, \the\NNODE\space nodes, root \the\DTROOT}%
        %% COLORMAP first: walls shade through it now, not only sprites, so
        %% it is no longer part of the thing loader and a map with no THINGS
        %% still gets its shading.  Textures follow, once every sidedef has
        %% had its say about which ones the map needs.
        \dtLoadColormap
        \dtLoadTextures
        \dtLoadFlats
      \fi
      %% ---- PLAYPAL: 768 bytes of RGB is exactly a PDF indexed colour space
      \dtFindLump{PLAYPAL}{0}%
      \ifnum\dtFOUND<0
        \message{[wad] no PLAYPAL -- keeping the built-in palette.}%
      \else
        \xdef\DTPAL{\pdffiledump offset \the\dtLPOS\space length 768
          {\dtWADFILE}}%
        \message{[wad] PLAYPAL installed (the real Doom palette)}%
      \fi
      %% ---- THINGS: player 1 start is thing type 1
      \dtFindLump{THINGS}{\the\dtMAP}%
      \ifnum\dtFOUND<0\else
        %% \tmpa cannot be the loop bound here: \dtnorm below writes to it,
        %% and so does \dtPointSide in the second pass.
        \dtNTH=\dtLSIZE \divide\dtNTH by10
        \dtJ=0 \dtFOUND=-1
        \dtWhile{\dtJ<\dtNTH}{%
          \ifnum\dtFOUND<0
            \ifnum\dtUu{\number\numexpr\dtLPOS+10*\dtJ+6\relax}=1
              \dtFOUND=\dtJ
              \PX=\numexpr\dtS{\number\numexpr\dtLPOS+10*\dtJ\relax}%
                   *65536\relax
              \PY=\numexpr\dtS{\number\numexpr\dtLPOS+10*\dtJ+2\relax}%
                   *65536\relax
              %% Doom stores the angle in degrees; 8192 units to a full turn
              \dtnorm\PA{\dtUu{\number\numexpr\dtLPOS+10*\dtJ+4\relax}%
                *8192/360}%
            \fi
          \fi
          \advance\dtJ by1\relax }%
        \ifnum\dtFOUND<0\else
          \message{[wad] player start
            \number\numexpr\PX/65536\relax,\number\numexpr\PY/65536\relax\space
            angle \the\PA}%
        \fi
        %% ---- every other thing, for the sprite pass.  A thing stands on the
        %% ---- floor of whatever sector contains it, which is exactly what
        %% ---- M1's point locator answers -- so this needs the BSP loaded.
        \ifnum\NNODE>0
          \dtJ=0
          \dtWhile{\dtJ<\dtNTH}{%
            \dtTT=\dtUu{\number\numexpr\dtLPOS+10*\dtJ+6\relax}\relax
            \dtTF=\dtUu{\number\numexpr\dtLPOS+10*\dtJ+8\relax}\relax
            \dtTB=\dtTF \divide\dtTB by16
            \tmpd=\dtTB \divide\tmpd by2 \multiply\tmpd by2
            \advance\dtTB by-\tmpd
            %% flags bit 4 is MTF_NOTSINGLE: multiplayer only
            \ifnum\dtTB=1 \else
              \dtPSX=\dtS{\number\numexpr\dtLPOS+10*\dtJ\relax}\relax
              \dtPSY=\dtS{\number\numexpr\dtLPOS+10*\dtJ+2\relax}\relax
              \dtPointInSSec
              \edef\dt@next{\noexpand\dtAddThing{\the\dtPSX}{\the\dtPSY}%
                {\the\dtTT}{\csname dtSf\the\dtSECN\endcsname}}%
              \dt@next
              %% the sector's light, for shading the sprite through COLORMAP,
              %% and the thing's own facing, for picking its rotation.  Doom
              %% stores the angle in degrees.  \dtAddThing has already
              %% advanced \NTHING.
              \expandafter\xdef\csname dtHl\number\numexpr\NTHING-1\relax
                \endcsname{\csname dtSl\the\dtSECN\endcsname}%
              %% ...and the sector itself, which the point locator has just
              %% answered and nothing else here ever asks it again: a
              %% teleporter finds its destination by it, and a thing in a
              %% strobing room is lit by it frame by frame.
              \expandafter\xdef\csname dtHs\number\numexpr\NTHING-1\relax
                \endcsname{\the\dtSECN}%
              \expandafter\xdef\csname dtHa\number\numexpr\NTHING-1\relax
                \endcsname{\number\numexpr
                  \dtUu{\number\numexpr\dtLPOS+10*\dtJ+4\relax}*8192/360\relax}%
            \fi
            \advance\dtJ by1\relax }%
          \message{[wad] \the\NTHING\space things placed}%
          %% Only the pictures this map actually needs: a full sprite set is
          %% 1381 lumps and a map uses a dozen.
          \dtJ=0
          \dtWhile{\dtJ<\NTHING}{%
            \edef\dt@wty{\csname dtHt\the\dtJ\endcsname}%
            \edef\dt@k{\dtSpriteKey{\dt@wty}}%
            \ifx\dt@k\dt@empty\else
              %% An animated type needs every frame of its sequence, not just
              %% the one \dtSN names, or \dt@animframe asks for a picture that
              %% was never read and the thing drops to its rotation 0.
              \ifcsname dtMA\dt@wty\endcsname
                \expandafter\dt@wantall\expandafter{\dt@k}%
              \else
                \edef\dt@next{\noexpand\dtWantSprite{\dt@k}}\dt@next
              \fi
              %% ...and a monster needs every frame of every STATE, which is
              %% the same argument one level up: walking, shooting, flinching
              %% and dying are four more sequences and a zombieman that dies
              %% into a picture nobody read turns into a coloured billboard
              %% on its last frame.  A type that is not in the monster table
              %% costs one \ifcsname that misses.
              %%
              %% Not under \DTNOMON, which is why this reads a configuration
              %% macro the way \dtLoadTextures reads \DTNOTEX: with nothing
              %% spawned, every monster stays on its spawn frame for the
              %% whole run and the other four sequences are ninety pictures
              %% nobody will look at.  It takes the doom2 MAP01 load back to
              %% what it was before M15, which is what the dozen fixtures
              %% that set \DTNOMON want anyway.
              \expandafter\ifx\csname DTNOMON\endcsname\relax
                \ifcsname dtOq\dt@wty\endcsname
                  \expandafter\dt@wantstates\expandafter{\dt@k}%
                \fi
              \fi
            \fi
            %% Outside the \ifx above that tests dt@k, because a missile's
            %% pictures do not hang off this thing's sprite key at all -- and
            %% outside \DTNOMON, because the player throws two of the three
            %% whether or not anything else on the map is awake.  Under
            %% \DTNOMIS none of it is read, which is the load M15 left.
            \expandafter\ifx\csname DTNOMIS\endcsname\relax
              \ifcsname dtJm\dt@wty\endcsname
                \edef\dt@next{\noexpand\dt@wantmissile
                  {\csname dtJm\dt@wty\endcsname}}\dt@next
              \fi
            \fi
            \advance\dtJ by1\relax }%
          %% The weapons are ordinary sprite lumps, so they ride along on the
          %% sprite pass -- which is what makes picking one up show it: the
          %% HUD reads dtUS<weapon> and the picture is already there.
          \dtWantSprite{PUNGA}\dtWantSprite{PISGA}\dtWantSprite{SHTGA}%
          \dtWantSprite{CHGGA}\dtWantSprite{MISGA}\dtWantSprite{PLSGA}%
          \dtWantSprite{BFGGA}\dtWantSprite{SAWGA}\dtWantSprite{SHT2A}%
          \dtLoadSprites
          %% The bar and its numbers are not sprites: they live in the main
          %% directory, so they are wanted by name and resolved in one pass.
          \dt@wanthud
          \dtLoadGfxAll
          \message{[wad] hud: \the\dtGN\space pictures}%
          \ifcsname dtPwSTBAR\endcsname
            \message{[wad] STBAR: \csname dtPwSTBAR\endcsname x\csname dtPhSTBAR\endcsname}%
          \else \message{[wad] no STBAR}\fi
        \fi
      \fi
    \fi
  \fi}
\catcode`\@=12
